\uuid{B7eu}
\exo7id{188}
\titre{exo7 188}
\auteur{bodin}
\organisation{exo7}
\datecreate{1998-09-01}
\isIndication{false}
\isCorrection{false}
\chapitre{Injection, surjection, bijection}
\sousChapitre{Injection, surjection}
\cpge{mpsi-mp2i.raisonnement-vocabulaire-ensembliste.c, pcsi.raisonnement-vocabulaire-ensembliste.d, ptsi.raisonnement-vocabulaire-ensembliste.d, mpsi-mp2i.fonctions-valeurs-reelles-complexes.b, pcsi.techniques-calcul-fonctions.b, ptsi.techniques-calcul-fonctions.b, mpsi-mp2i.limites-continuite.c, pcsi.limites-continuite.c, ptsi.limites-continuite.c, bcpst1.vocabulaire-applications, bcpst1.fonctions-reelles-usuelles.b}
\module{Algèbre}
\niveau{L1}
\difficulte{}

\contenu{
\question{
Soit $f  : \Rr \rightarrow \Rr$ d\'efinie par $f(x) = x^3-x$.\\
$f$ est-elle injective ? surjective ? D\'eterminer $f^{-1}([-1,1])$ et $f(\Rr_+)$.
}
}
